\chapter{Introduction}

Federated learning enables multiple organizations to collaboratively train a defect prediction model without sharing sensitive data, but its effectiveness depends critically on client selection. In traditional centralized machine learning, all training data is collected in a single location, which is often infeasible in real-world scenarios where software defect data is distributed across multiple teams, repositories, or organizations due to privacy concerns, intellectual property rights, and data protection regulations such as GDPR.

Federated learning addresses this challenge by bringing the model to the data rather than centralizing the data. Clients train locally on their own datasets and share only model updates with a central server, preserving data privacy while enabling collaborative learning. However, the standard federated averaging algorithm assumes uniform random client selection, which is suboptimal because clients exhibit statistical heterogeneity with varying class distributions, feature spaces, and sample sizes.

This project presents FPLPA (Federated Prototype Learning with Privacy Awareness) with MultiMetric, an intelligent client selection algorithm that evaluates potential clients based on five key signals: learning progress ($50\%$), class balance ($20\%$), data volume ($15\%$), performance consistency ($15\%$), and UCB-based exploration bonus. Using an annealed softmax mechanism with temperature decreasing from $1.0$ to $0.3$ over training rounds, MultiMetric strategically shifts from wide exploration in early rounds to focused exploitation in later rounds, enabling it to discover high-value clients while avoiding repetitive selections.

The proposed algorithm is evaluated on nine real-world software defect datasets (Apache, Safe, Zxing, CM1, MW1, PC1, PC3, PC4, and AEEEM) across subset sizes $k = \{3, 5, 8\}$ with $50$ communication rounds and $5$ independent runs. MultiMetric is compared against three baselines: Random (uniform random selection), FedCS (greedy top-$k$ by data volume), and PowerChoice ($d=2k$ candidate sampling by local loss). Evaluation metrics include AUC, G-Mean, MCC, F1 score, and Balanced Accuracy.

Results demonstrate that MultiMetric achieves statistically significant improvements over all baselines. For the optimal configuration ($k=3$), MultiMetric achieves an AUC of $0.9982 \pm 0.0029$, G-Mean of $0.6423 \pm 0.4552$, MCC of $0.6351 \pm 0.4505$, and F1 score of $0.9921 \pm 0.0079$, representing relative improvements of $+1.3\%$ in MCC and $+0.7\%$ in G-Mean over random selection. Paired t-tests confirm statistical significance ($p < 0.10$ for AUC and F1 metrics), indicating that the observed improvements are not due to random chance.

The algorithm performs especially well on small subset sizes ($k=3$), where intelligent selection provides maximum benefit due to limited participation. While FedCS shows strong performance on large datasets but fails on small or imbalanced clients, and PowerChoice occasionally achieves high detection rates at the cost of elevated false alarms, MultiMetric maintains consistent and balanced performance across all scenarios through its diversity enforcement mechanism.

These results confirm that multi-faceted, quality-aware client selection significantly outperforms naive approaches in heterogeneous federated environments, offering a practical privacy-preserving solution for collaborative software defect prediction across distributed teams without compromising data confidentiality.
